\section{Conclusions}
\label{sec:conclusions}

We presented a reproducible benchmark for high-resolution (1\,m) snow-depth
prediction in the Izas experimental catchment of the Central Spanish Pyrenees,
built from a 27-date airborne LiDAR time series and a rich stack of
topographic, wind-sheltering and snow-history predictors provided in
collaboration with IPE--CSIC. Using a strictly temporal train/validation/test
split that forces extrapolation to an exceptionally snowy year, we compared a
Random Forest baseline with U-Net and ResUNet++ convolutional networks, and we
evaluated all models with pattern-aware spatial metrics (\SPAEF{}, \MSPAEF{})
in addition to conventional pixel-wise scores.

Our main conclusions are:

\begin{itemize}
  \item \textbf{ResUNet++ with a spatial loss is the best and most reliable
  model.} At the optimal spatial-loss weight $\lpear=0.4$ it achieves
  \Rsq{}~$=0.208\pm0.027$ and \SPAEF{}~$=0.297\pm0.018$, outperforming both the
  Random Forest (\SPAEF{}~$=0.107$) and a far less stable plain U-Net
  (\Rsq{}~$=0.054\pm0.124$, \SPAEF{}~$=0.186$).

  \item \textbf{Pattern-aware evaluation is essential.} \SPAEF{} reverses the
  ranking of the Random Forest and the U-Net relative to \Rsq{}, and reveals
  that the per-pixel RF reproduces the spatial pattern almost three times worse
  than ResUNet++ -- effects that an \Rsq{}-only assessment would miss. We
  recommend reporting \SPAEF{}/\MSPAEF{} and multi-seed statistics as standard
  practice in high-resolution snow mapping.

  \item \textbf{The hybrid spatial loss helps twice.} Blending pixel-wise MSE
  with a spatial-correlation term improves both accuracy and pattern fidelity
  and substantially reduces the variance across seeds, with a clear optimum at
  $\lpear=0.4$.

  \item \textbf{Snow history and meso-scale topography are the key
  predictors,} whereas instantaneous satellite snow-cover extent is nearly
  redundant, and scalar meteorological forcing does not significantly improve
  the prediction for this catchment.
\end{itemize}

By coupling a stable, attention-augmented architecture with a pattern-oriented
loss and evaluation, this work provides both a methodological template and a
public benchmark for high-resolution snow-depth mapping. Future work will
extend the approach to multiple catchments and snow climates, incorporate
spatially distributed meteorological forcing and cloud-aware snow-history
inputs, and assess transferability across years with contrasting snow
conditions.
